\clearpage
\section{Implementation update: 14 November 2026}
\label{sec:implementation-update}

This section records implementation after the mid-term experiment. The bank now contains
135 team-checked questions and 95 misconceptions. Code, schemas, controls and evaluation
support have changed; the historical B1/A2 results above are not results for these successors.

\subsection{Evidence-backed candidate-belief modelling}
The historical A2 confirms a misconception from two replies, which can belong to the same
question. A2Q instead counts distinct main-question identifiers: repeated supporting replies
within one question supply provenance but not a second independent question. Confirmation
across questions is still model-generated evidence, not independent human validation.

A2E extends A2Q with an evidence prompt. Each positive vote must name an allowed misconception,
give a non-empty exact substring of the current candidate reply, and explain its relevance.
Code validates the identifier, quote and offsets. Unsupported positive votes become gaps.
A positive diagnosis needs a strict majority of the configured diagnosis votes; missing or
unusable votes do not lower that majority requirement. Matching quotations establish
provenance, while the semantic relationship between quote and misconception still depends on
the model.

The tracker stores question, turn, reply index, reply text, vote share and supporting evidence.
Correct main replies are also diagnosed in A2E, without adding follow-ups to those replies,
so explicit corrections may contradict previously attributed beliefs. Accepted contradictory
evidence marks an existing finding \emph{disputed}, with confirmation false, for the remainder
of that interview. Silence or an unrelated correct answer is not itself a contradiction.
This increases diagnosis calls and does not establish improved precision or shorter interviews.

A4 is unchanged: one main-answer diagnosis, generic follow-ups, and one-reply confirmation.
The A2E/A4 comparison tests this combined package. The optional tracker-reset control A2M
is not implemented, so there is no isolated experimental claim about persistent memory.

\subsection{Level uncertainty and stopping alternative}
New reports include the posterior mass in each level band, and the mass in the reported
EAP level's band. The bands retain cuts $-1.25$, $-0.25$, $0.75$ on the existing finite
ability grid; each cut belongs to its upper band. These are probabilities conditional on the
score model, not calibrated probabilities of the simulated ground-truth labels.

A2EB differs from A2E only in its stopping rule: every skill requires at least two main
answers, SE below 0.6, and mass at least 0.7 in its reported level. A2B is the matched
alternative for historical A2. The 15-main-question cap and the fit/review rules are unchanged.
Neither alternative is adopted as the main condition.

On the separate 12-candidate stub validation set, A2E reported 30/36 skill levels correctly
and A2EB 31/36, with mean exchanges 15.58 and 17.67 respectively. Every A2EB interview hit
the cap. Its paired accuracy difference was 2.8 points (95\% candidate-bootstrap interval
$-5.6$ to $+11.1$ points); true-level-0 errors remained 1/2 assessed skills for both.
This does not establish an improvement. Eliminating confident stops trivially eliminates
errors conditional on confident stops; all-stop error and cap frequencies remain necessary.

\subsection{Tool calling and trusted SQL verification}
The existing agent driver is now exposed by a runnable demonstration. It calls rule-enforcing
question, grading, follow-up, scorecard and stopping tools. Deliberately invalid actions are
refused. A summary-consistency check compares the written conclusion with the computed outcome;
it is a heuristic rather than a proof of faithful language generation. The main evaluation
still uses the code driver; live agent overhead and quality are not established by a script.

A2EV is the optional A2E condition with SQL verification. For eligible bank SQL read-code
questions, the tool executes the bank's query on its supplied tables using in-memory SQLite.
Only outputs matching the bank's checked expected result are supplied to the grader.
Mismatch or execution error supplies no verified evidence and flags the report. Candidate
text is never executed, candidate prompts are unchanged, and result rows are not given to the
agent. All 18 eligible bank queries pass the offline execution checks. This validates bank
execution, not the grader's semantic judgements or arbitrary candidate-written SQL.

The reproducible offline example uses VAL-P-004. A2E and A2EV each produce 14 exchanges;
their obedient code and scripted-agent runs agree on questions, skill levels and decision.
The probed agents take 56 and 61 scripted steps respectively, with consistent summaries.
These steps make no network calls and do not measure live model performance or latency.

\subsection{Evaluation and stopped live smoke}
ADAPT12 is an actual preset with adaptive selection, 12 main questions, no diagnosis and no
follow-ups: exactly 12 exchanges. Its selection retains weak-topic ranking as part of the
adaptive package. A5E differs from A2E only in job-blind retrieval. A manifest fingerprints
source, bank, prompts, setup, configurations and candidate files; CRLF is normalised to LF for
cross-platform hashing. Incompatible runs cannot be silently resumed. Pairing, failures,
exclusions and denominators are recorded, and human-supported endpoints remain unmeasured.

The attempted all-Groq pilot used Qwen 3.8 27B for grading and GPT-OSS 120B for simulation.
Two initial B1 attempts failed because Groq rejected generated JSON. A narrow client fix now
classifies that exact output-error code like malformed HTTP-200 JSON, allowing the existing
vote-handling code to reject unusable votes. Invalid requests remain fatal; rejected output
is not cached or accepted, and an all-invalid batch still fails. Fewer usable grading votes
can increase agreement among the survivors, so vote counts require review.

Under a fresh manifest, one B1 interview completed: 12 main questions, no interview errors,
and 37 successful interviewer completions, of which 27 were uncached. It took 11 minutes
24 seconds, including waits, simulation and cache reuse. The A2E smoke then failed with no
usable diagnosis votes. The full 60-interview pilot never started and no process remains active.
These observations establish a compatibility/resource problem, not an accuracy comparison.

The successful-completion meter excludes failed HTTP attempts and simulator usage; token
totals can include locally cached completions. It must not be treated as a complete invoice.
The independent human checks were deferred. No revised-bank live improvement, isolated memory
benefit, or SQL-assisted grading gain is claimed. Final comparative results remain pending.
